\begin{tikzpicture}[font=\small,>=Stealth,
  box/.style={draw=Primary,fill=PrimaryLight,rounded corners=2pt,minimum width=2.9cm,minimum height=1.05cm,align=center,font=\scriptsize,text=PrimaryDark}]
  \node[box] (p) at (-4.6,0) {降雨情景参数\\雨量·历时·分布};
  \node[box] (m) at (-1.2,0) {产汇流与\\淹没计算};
  \node[draw=Primary!40,rounded corners=3pt,minimum width=4.6cm,minimum height=3.2cm] (v) at (3.4,0) {};
  \node[anchor=north,font=\scriptsize,text=PrimaryDark] at ($(v.north)+(0,-.12)$) {淹没范围与相对水深（示意）};
  \begin{scope}
    \clip (1.2,-1.45) rectangle (5.6,1.25);
    \fill[Primary!15] (3.4,-.3) ellipse (1.9 and .75);
    \fill[Primary!35] (3.65,-.3) ellipse (1.2 and .5);
    \fill[Primary!60] (4.1,-.25) ellipse (.5 and .25);
  \end{scope}
  \node[font=\scriptsize,text=TextGray] at (3.4,-1.3) {同一时刻：深色表示相对水深较大};
  \draw[->,thick,Primary] (p)--(m);
  \draw[->,thick,Primary] (m)--(v);
  \draw[->,thick,Accent,dashed] (v.south)--(3.4,-2.25)
    --node[below,font=\scriptsize]{调整情景后重新计算} (-4.6,-2.25)--(p.south);
\end{tikzpicture}
